% =============================================================
% Annexure C: DNN Architecture and Training Procedure
% =============================================================
\section{DNN Architecture and Training Procedure}
\label{annexure:dnn_training}

The FunnelDNN model uses a feed-forward architecture with monotonically decreasing hidden layer widths: $22 \to 256 \to 128 \to 64 \to 32 \to 16 \to 1$. All layers are bias-free, and the CELU activation function is applied after each hidden layer. The output layer is linear with no activation.

% -------------------------------------------------------------
\subsection{Training Procedure}
\label{annexure:train_model}

Training minimises the mean squared error loss on normalised data using the NAdam optimiser. Validation-based early stopping monitors the dot-product $R^2$ on the validation set every 5 epochs, with a patience of 10 checks.

\begin{algorithm}[htbp]
\caption{FunnelDNN training with early stopping.}
\label{alg:dnn_training}
\begin{algorithmic}[1]
\Require Standardised data $(\mathbf{X}_{\mathrm{train}}, \mathbf{y}_{\mathrm{train}}, \mathbf{X}_{\mathrm{val}}, \mathbf{y}_{\mathrm{val}})$, learning rate $\eta$, batch size $B$, patience $P$, validation interval $V$
\Ensure Trained model state $\boldsymbol{\theta}^*$
\State Initialise FunnelDNN with seed 42
\State $\boldsymbol{\theta}^* \gets \boldsymbol{\theta}_0$, \; $R^{2*} \gets -\infty$, \; $c \gets 0$
\For{$t = 1, 2, \ldots, T_{\max}$}
    \For{each mini-batch $(\mathbf{X}_b, \mathbf{y}_b)$ of size $B$}
        \State Compute $\mathcal{L} = \mathrm{MSE}\bigl(f_{\boldsymbol{\theta}}(\mathbf{X}_b),\, \mathbf{y}_b\bigr)$
        \State Update $\boldsymbol{\theta}$ via NAdam with learning rate $\eta$
    \EndFor
    \If{$t \bmod V = 0$}
        \State Compute $R^2_{\mathrm{val}} = r^2_{\mathrm{dot}}\bigl(\mathbf{y}_{\mathrm{val}},\, f_{\boldsymbol{\theta}}(\mathbf{X}_{\mathrm{val}})\bigr)$
        \If{$R^2_{\mathrm{val}} > R^{2*}$}
            \State $R^{2*} \gets R^2_{\mathrm{val}}$, \; $\boldsymbol{\theta}^* \gets \boldsymbol{\theta}$, \; $c \gets 0$
        \Else
            \State $c \gets c + 1$
        \EndIf
        \If{$c \geq P$}
            \State \textbf{break} \Comment{Early stopping triggered}
        \EndIf
    \EndIf
\EndFor
\State \Return $\boldsymbol{\theta}^*$
\end{algorithmic}
\end{algorithm}

% -------------------------------------------------------------
\subsection{Evaluation Metrics}
\label{annexure:eval_helpers}

Three metrics are computed on the original Rand scale after de-normalisation:
\begin{itemize}
    \item $R^2 = 1 - \mathrm{SS}_{\mathrm{res}} / \mathrm{SS}_{\mathrm{tot}}$, where $\mathrm{SS}_{\mathrm{res}} = \sum_i (y_i - \hat{y}_i)^2$ and $\mathrm{SS}_{\mathrm{tot}} = \sum_i (y_i - \bar{y})^2$.
    \item $\mathrm{RMSE} = \sqrt{n^{-1} \sum_i (y_i - \hat{y}_i)^2}$.
    \item $\mathrm{MAE} = n^{-1} \sum_i |y_i - \hat{y}_i|$.
\end{itemize}

During training, the dot-product $R^2$ (squared correlation between predicted and observed values) is used on normalised data for computational speed.

% -------------------------------------------------------------
\subsection{Checkpoint Contents}
\label{annexure:final_model}

The final trained model is saved as \texttt{chapter3\_final\_model.pth}, a PyTorch checkpoint containing:
\begin{itemize}
    \item Model weight tensors (state dictionary).
    \item Architecture specification (layer widths, activation, bias setting).
    \item Selected hyperparameters (learning rate, batch size, optimiser, momentum).
    \item Standardisation parameters ($\boldsymbol{\mu}_X$, $\boldsymbol{\sigma}_X$, $\mu_y$, $\sigma_y$).
    \item Feature names and evaluation metrics.
\end{itemize}
